\scanpage{57}
\mspage{53}
% archive sheet number: 56
qu'on a un hom injectif, qui identifie $\unsure{\mathcal A|\mathscr A}$ à un sous-atelier de $\mathrm{Atspat}(\Sigma)$, \deleted{\unsure{soit}} \inserted{noté \unsure{encore}} $\deleted{\unsure{\mathcal B|B}}\unsure{\mathcal A|A}$

Ainsi, on trouve, associé à $\mathcal A$, \unsure{$x$|:}

$1^{\circ}$) un Spatios $\Sigma$

$2^{\circ}$) un sous-atelier $\deleted{\unsure{\mathcal B|B}}\mathcal A\subset\mathrm{At}\,\mathrm{spat}(\Sigma)$. \deleted{\unsure{i.e.}}

Sous-atelier signifie ici, simplement, une \uline{partie}\inserted{de $\mathrm{Atspat}$}, qui \unsure{satisfasse|satisfait} la condition Sous at (cf p. 27) --- mais en fait, c'est \unsure{$=$|même} une partie \uline{fermée pour} $\trianglelefteq$.
Cette condition joue le rôle de la condition $\deleted{\unsure{\mathrm O|\mathrm{Ob}}}\mathrm{Atens}\ 1$ (cf \unsure{CF|CK} $\overline{\text{\uline{VIII}}}$\inserted{ p.~\unsure{309|39}}) pour les ateliers ensemblistes, rapportés à $L$\inserted{ (ens. des \unsure{lieux})} comme ens. de base \unsure{\}|?}. Il faut une condition qui $\underset{\text{-- -- -- --}}{\text{\deleted{\unsure{tienne}}}}$ lieu de $\mathrm{Atens}\ 2$ ($\forall x\in L,\ \{\{x\}\}\in A$). Ce sera Atspat 2. Ci dessous.

\deleted{Donc} On va donc rapporter la structure sur l'ens de base $\Sigma$, avec comme données $(\leq,\deleted{\unsure{\mathrel{|o|}|\trianglelefteq}}C,\mathcal A)$\deleted{\illegible{1}}, où $\leq,C$ font de $\Sigma$ un \unsure{spatios|Spatios}, et où $\mathcal A$ est un \deleted{\unsure{partie|part}} ens. \unsure{des|de} figures spatiales élémentaires \deleted{$\unsure{\mathcal A|A}$} relativement à $\Sigma$
\[
\mathcal A\subset\mathrm{At}\,\mathrm{spat}(\Sigma),
\]
